\subsection*{CU-07: Vincular la cuenta de invitado}
\addcontentsline{toc}{subsection}{CU-07: Vincular la cuenta de invitado}
\label{cu-07}

\begin{fichacu}{CU-07}{Vincular la cuenta de invitado}
  \crudFila{Responsable}{José Eduardo Prior Hernández}
  \crudFila{Actor principal}{Invitado}
  \crudFila{Actores de sistema}{Servidor de partidas, Servicio de correo, Base de datos}
  \crudFila{Prototipos}{9e, 6d}
  \crudFila{Objetivo}{Convertir una cuenta de invitado en una cuenta registrada conservando íntegro su progreso: elo, monedas, objetos, partidas e historial.}
  \crudFila{Disparador}{El Invitado selecciona ``Crear cuenta'' desde el aviso permanente de la \texttt{GUI\_\allowbreak{}MainMenu} o desde la \texttt{GUI\_\allowbreak{}AccountSettings}.}
  \textbf{Precondiciones} & \cuCelda{\begin{cureglas}
      \item \textbf{\mbox{PRE-01}} El Invitado tiene una \textit{Sesion} abierta sobre un \textit{Usuario} con \texttt{tipo\_\allowbreak{}cuenta} en \texttt{INVITADO}.
      \item \textbf{\mbox{PRE-02}} El Servidor de partidas está en ejecución y tiene conexión disponible con la Base de datos.
      \item \textbf{\mbox{PRE-03}} Existe una versión vigente de los términos de uso y su texto está disponible en el idioma de la interfaz.
    \end{cureglas}} \\ \hline
  \textbf{Flujo normal} & \cuCelda{\begin{cuenum}[start=1]
      \item El SJM despliega la \texttt{GUI\_\allowbreak{}LinkAccount}, solicitando el \texttt{nickname} definitivo, el \texttt{correo}, la contraseña y su confirmación y la \texttt{fecha\_\allowbreak{}nacimiento}, junto con la casilla de aceptación de los términos de uso. El formulario indica de forma visible qué progreso se conserva. \textit{(Ver \mbox{FA-01})}
      \item El Invitado captura los datos, marca la casilla y da click en ``Crear cuenta''.
      \item El SJM valida que todos los campos obligatorios estén completos y con el formato esperado. \textit{(Ver \mbox{FA-02})}
      \item El SJM valida que la contraseña cumpla la política de \mbox{RN-03} y que coincida con su confirmación. \textit{(Ver \mbox{FA-03})}
      \item El SJM valida que la casilla de los términos esté marcada. \textit{(Ver \mbox{FA-04})}
      \item El SJM calcula la edad a partir de la \texttt{fecha\_\allowbreak{}nacimiento} y verifica que sea de ocho años o más. \textit{(Ver \mbox{FA-05})}
      \item El SJM envía el mensaje \texttt{LINK\_\allowbreak{}ACCOUNT\_\allowbreak{}REQUEST} con los datos capturados, la contraseña, la versión de los términos aceptada, el idioma en que se mostró su texto y el token de la sesión de invitado. \textit{(Ver \mbox{EX-05}, \mbox{EX-06}, \mbox{EX-02})}
      \item El Servidor de partidas verifica que el token corresponda a una \textit{Sesion} abierta de un \textit{Usuario} con \texttt{tipo\_\allowbreak{}cuenta} en \texttt{INVITADO}. \textit{(Ver \mbox{FA-06})}
    \end{cuenum}} \\
   & \cuCelda{\begin{cuenum}[start=9]
      \item El Servidor de partidas vuelve a aplicar sobre los datos recibidos todas las validaciones de los pasos 3 a 6, y verifica que el idioma de los términos sea uno de los admitidos, \texttt{es-MX} o \texttt{en} (\mbox{D-21}), sin dar por buena la validación del cliente (\mbox{RN-09}). \textit{(Ver \mbox{FA-07})}
      \item El Servidor de partidas verifica que el \texttt{nickname} no contenga términos del catálogo de \textit{PalabraProhibida}. \textit{(Ver \mbox{FA-08})}
      \item El Servidor de partidas verifica que el \texttt{nickname} no esté en uso por otro \textit{Usuario}. \textit{(Ver \mbox{FA-09})}
      \item El Servidor de partidas normaliza el \texttt{correo} a minúsculas y verifica que no esté registrado por otro \textit{Usuario}. \textit{(Ver \mbox{FA-10})}
      \item El Servidor de partidas verifica que la versión de los términos recibida sea la vigente. \textit{(Ver \mbox{FA-11})}
      \item El Servidor de partidas inicia una transacción sobre la Base de datos.
      \item El Servidor de partidas genera la \texttt{contrasena\_\allowbreak{}sal} y deriva la \texttt{contrasena\_\allowbreak{}hash} de la contraseña recibida (\mbox{RN-02}). \textit{(Ver \mbox{EX-03})}
      \item El Servidor de partidas actualiza el \textit{Usuario} existente: cambia \texttt{tipo\_\allowbreak{}cuenta} a \texttt{REGISTRADA}, \texttt{estado\_\allowbreak{}cuenta} a \texttt{PENDIENTE}, y escribe el \texttt{nickname} definitivo, el \texttt{correo}, la \texttt{contrasena\_\allowbreak{}hash}, la \texttt{contrasena\_\allowbreak{}sal}, la \texttt{fecha\_\allowbreak{}nacimiento} y un \texttt{codigo\_\allowbreak{}amigo} único, que la cuenta de invitado no tenía porque no podía tener amigos (\mbox{RN-01}). \textit{(Ver \mbox{EX-03})}
    \end{cuenum}} \\
   & \cuCelda{\begin{cuenum}[start=17]
      \item El Servidor de partidas registra la \textit{AceptacionTerminos} con la versión aceptada, el idioma en que se mostró su texto (\mbox{D-21}), la fecha y la dirección de red de origen. \textit{(Ver \mbox{EX-03})}
      \item El Servidor de partidas genera un \textit{TokenVerificacionCorreo} para el \textit{Usuario}, con su fecha de expiración. \textit{(Ver \mbox{EX-03})}
      \item El Servidor de partidas confirma la transacción. \textit{(Ver \mbox{EX-04})}
      \item El Servidor de partidas solicita al Servicio de correo el envío del mensaje de verificación, en el \texttt{idioma\_\allowbreak{}preferido} del \textit{Usuario} (\mbox{RN-07}). \textit{(Ver \mbox{EX-07})}
      \item El Servidor de partidas actualiza \texttt{fecha\_\allowbreak{}ultimo\_\allowbreak{}envio\_\allowbreak{}verificacion} del \textit{Usuario}.
      \item El Servidor de partidas registra la vinculación en la \textit{BitacoraAcceso} con el resultado \texttt{VINCULACION\_\allowbreak{}EXITOSA}.
      \item El Servidor de partidas cierra la \textit{Sesion} de invitado con \texttt{motivo\_\allowbreak{}cierre} igual a \texttt{CAMBIO\_\allowbreak{}CREDENCIALES} y responde con \texttt{LINK\_\allowbreak{}ACCOUNT\_\allowbreak{}OK} (\mbox{RN-08}).
      \item El SJM despliega la \texttt{GUI\_\allowbreak{}RegistrationSuccess}, indicando que se envió un correo de verificación a la dirección capturada, mostrada de forma parcialmente oculta, y que debe verificarla para volver a entrar.
    \end{cuenum}} \\
   & \cuCelda{\begin{cuenum}[start=25]
      \item El SJM descarta el token de invitado guardado en el dispositivo y despliega la \texttt{GUI\_\allowbreak{}Login}.
      \item Fin del caso de uso.
    \end{cuenum}} \\ \hline
  \textbf{Flujos alternos} & \cuCelda{\textbf{\mbox{FA-01}: El Invitado cancela la vinculación}\par
    \begin{cuenum}[start=1]
      \item El Invitado selecciona ``Cancelar''.
      \item El SJM descarta los datos capturados sin enviarlos y vuelve a la \texttt{GUI\_\allowbreak{}MainMenu}, con la sesión de invitado intacta.
      \item Fin del caso de uso.
    \end{cuenum}} \\
   & \cuCelda{\textbf{\mbox{FA-02}: Campos incompletos o con formato inválido}\par
    \begin{cuenum}[start=1]
      \item El SJM muestra la \texttt{GUI\_\allowbreak{}MessageWarning} indicando qué campos están incompletos o mal formados, los resalta y escribe junto a cada uno qué se espera de él.
      \item El flujo continúa en el paso 2 del FN. La validación es local y no llega al Servidor de partidas.
    \end{cuenum}} \\
   & \cuCelda{\textbf{\mbox{FA-03}: La contraseña no cumple la política o no coincide}\par
    \begin{cuenum}[start=1]
      \item El SJM muestra la \texttt{GUI\_\allowbreak{}MessageWarning} explicando en una frase qué le falta conforme a \mbox{RN-03}, o que las dos capturas no coinciden, y limpia ambos campos.
      \item El flujo continúa en el paso 2 del FN.
    \end{cuenum}} \\
   & \cuCelda{\textbf{\mbox{FA-04}: No se aceptaron los términos de uso}\par
    \begin{cuenum}[start=1]
      \item El SJM resalta la casilla de aceptación y escribe junto a ella que debe marcarse para continuar, sin abrir ninguna ventana adicional.
      \item El flujo continúa en el paso 2 del FN.
    \end{cuenum}} \\
   & \cuCelda{\textbf{\mbox{FA-05}: El Invitado no alcanza la edad mínima}\par
    \begin{cuenum}[start=1]
      \item El SJM muestra la \texttt{GUI\_\allowbreak{}MessageWarning} indicando que el juego está dirigido a personas de ocho años o más.
      \item El SJM descarta los datos capturados sin enviarlos al Servidor de partidas y sin escribir la \texttt{fecha\_\allowbreak{}nacimiento} en ninguna bitácora (\mbox{RN-10}).
      \item El SJM vuelve a la \texttt{GUI\_\allowbreak{}MainMenu}. La cuenta de invitado sigue funcionando: la edad mínima condiciona el registro, no el juego.
      \item Fin del caso de uso.
    \end{cuenum}} \\
   & \cuCelda{\textbf{\mbox{FA-06}: La sesión de invitado ya no es válida}\par
    \begin{cuenum}[start=1]
      \item El Servidor de partidas no encuentra una \textit{Sesion} abierta de invitado para el token recibido: pudo cerrarse desde otro dispositivo, expirar o haber sido ya vinculada.
      \item El Servidor de partidas responde con \texttt{SESSION\_\allowbreak{}INVALID}.
      \item El SJM muestra la \texttt{GUI\_\allowbreak{}MessageWarning} indicando que la sesión terminó y despliega la \texttt{GUI\_\allowbreak{}Login}.
      \item Fin del caso de uso.
    \end{cuenum}} \\
   & \cuCelda{\textbf{\mbox{FA-07}: El servidor rechaza datos que el cliente había aceptado}\par
    \begin{cuenum}[start=1]
      \item El Servidor de partidas registra en la bitácora del servidor la discrepancia entre lo que el cliente dio por válido y lo que él rechaza, junto con la \texttt{version\_\allowbreak{}cliente} y la dirección de red de origen, por tratarse de un indicio de cliente modificado.
      \item El Servidor de partidas responde con \texttt{LINK\_\allowbreak{}ACCOUNT\_\allowbreak{}FAILED} y el motivo \texttt{DATOS\_\allowbreak{}INVALIDOS}, enumerando los campos rechazados.
      \item El SJM muestra la \texttt{GUI\_\allowbreak{}MessageWarning} con los campos rechazados y los resalta.
      \item El flujo continúa en el paso 2 del FN.
    \end{cuenum}} \\
   & \cuCelda{\textbf{\mbox{FA-08}: El nickname contiene lenguaje inapropiado}\par
    \begin{cuenum}[start=1]
      \item El Servidor de partidas responde con \texttt{LINK\_\allowbreak{}ACCOUNT\_\allowbreak{}FAILED} y el motivo \texttt{NICKNAME\_\allowbreak{}NO\_\allowbreak{}PERMITIDO}, sin indicar qué término activó el filtro, para no ir revelando el catálogo a base de intentos.
      \item El Servidor de partidas registra el rechazo en la bitácora del servidor.
      \item El SJM muestra la \texttt{GUI\_\allowbreak{}MessageWarning} indicando que ese \texttt{nickname} no puede usarse y limpia el campo.
      \item El flujo continúa en el paso 2 del FN.
    \end{cuenum}} \\
   & \cuCelda{\textbf{\mbox{FA-09}: El nickname ya está en uso}\par
    \begin{cuenum}[start=1]
      \item El Servidor de partidas arma hasta tres sugerencias de \texttt{nickname} disponibles a partir del capturado.
      \item El Servidor de partidas responde con \texttt{LINK\_\allowbreak{}ACCOUNT\_\allowbreak{}FAILED}, el motivo \texttt{NICKNAME\_\allowbreak{}EN\_\allowbreak{}USO} y las sugerencias.
      \item El SJM muestra la \texttt{GUI\_\allowbreak{}MessageWarning} con las sugerencias como opciones seleccionables.
      \item El flujo continúa en el paso 2 del FN.
    \end{cuenum}} \\
   & \cuCelda{\textbf{\mbox{FA-10}: El correo ya está registrado}\par
    \begin{cuenum}[start=1]
      \item El Servidor de partidas responde con \texttt{LINK\_\allowbreak{}ACCOUNT\_\allowbreak{}FAILED} y el motivo \texttt{CORREO\_\allowbreak{}EN\_\allowbreak{}USO}.
      \item El SJM muestra la \texttt{GUI\_\allowbreak{}MessageWarning} advirtiendo que ya existe una cuenta con ese correo y que vincular el invitado a esa cuenta no es posible: el progreso de dos cuentas no puede fundirse (\mbox{RN-05}).
      \item El SJM ofrece las opciones ``Usar otro correo'' y ``Aceptar''. Con la primera, el flujo continúa en el paso 2 del FN; con la segunda, termina el caso de uso con la cuenta de invitado intacta.
    \end{cuenum}} \\
   & \cuCelda{\textbf{\mbox{FA-11}: Los términos aceptados no son los vigentes}\par
    \begin{cuenum}[start=1]
      \item El Servidor de partidas responde con \texttt{LINK\_\allowbreak{}ACCOUNT\_\allowbreak{}FAILED}, el motivo \texttt{TERMINOS\_\allowbreak{}DESACTUALIZADOS} y el identificador de la versión vigente.
      \item El SJM carga el texto de la versión vigente, desmarca la casilla de aceptación y muestra la \texttt{GUI\_\allowbreak{}MessageWarning} indicando que los términos cambiaron mientras llenaba el formulario.
      \item El flujo continúa en el paso 2 del FN.
    \end{cuenum}} \\ \hline
  \textbf{Excepciones} & \cuCelda{\textbf{\mbox{EX-01}: No se puede establecer la conexión con el servidor}\par
    \begin{cuenum}[start=1]
      \item El SJM detecta que la conexión TCP no pudo establecerse dentro del tiempo límite y registra el fallo en la bitácora local.
      \item El SJM muestra la \texttt{GUI\_\allowbreak{}MessageError} con las opciones ``Reintentar'' y ``Aceptar''. Los datos capturados se conservan.
      \item Si el Invitado selecciona ``Reintentar'', el flujo continúa en el paso 7 del FN. Si selecciona ``Aceptar'', continúa en el paso 2.
    \end{cuenum}} \\
   & \cuCelda{\textbf{\mbox{EX-02}: La versión del protocolo es incompatible}\par
    \begin{cuenum}[start=1]
      \item El Servidor de partidas responde con \texttt{PROTOCOL\_\allowbreak{}VERSION\_\allowbreak{}MISMATCH} y cierra la conexión.
      \item El SJM muestra la \texttt{GUI\_\allowbreak{}MessageError} indicando que debe actualizarse.
      \item Fin del caso de uso.
    \end{cuenum}} \\
   & \cuCelda{\textbf{\mbox{EX-03}: Fallo al escribir en la base de datos}\par
    \begin{cuenum}[start=1]
      \item El Servidor de partidas revierte la transacción completa, de modo que el \textit{Usuario} siga siendo un invitado íntegro y no quede a medio convertir, con credenciales pero sin aceptación de términos o sin token de verificación (\mbox{RN-06}).
      \item El Servidor de partidas registra la excepción en la bitácora del servidor con un identificador de incidencia, sin incluir la contraseña recibida.
      \item El Servidor de partidas responde con \texttt{SERVER\_\allowbreak{}ERROR} y el identificador, y cierra la conexión.
      \item El SJM muestra la \texttt{GUI\_\allowbreak{}MessageError} con el identificador. Los datos capturados se conservan, salvo la contraseña.
      \item El flujo continúa en el paso 2 del FN.
    \end{cuenum}} \\
   & \cuCelda{\textbf{\mbox{EX-04}: Fallo al confirmar la transacción}\par
    \begin{cuenum}[start=1]
      \item El Servidor de partidas no obtiene confirmación de que la transacción se haya escrito y la trata como fallida.
      \item El Servidor de partidas registra la excepción señalando que el estado de la vinculación es indeterminado y responde con \texttt{SERVER\_\allowbreak{}ERROR}.
      \item El SJM muestra la \texttt{GUI\_\allowbreak{}MessageError} indicando que revise su correo antes de reintentar: si la vinculación sí se completó, un segundo intento terminará en el \mbox{FA-10}.
      \item Fin del caso de uso.
    \end{cuenum}} \\
   & \cuCelda{\textbf{\mbox{EX-05}: Se agota el tiempo de espera de la respuesta}\par
    \begin{cuenum}[start=1]
      \item El SJM detecta que transcurrió el tiempo límite sin recibir un mensaje completo.
      \item El SJM cierra la conexión TCP y descarta el intercambio, sin suponer que el servidor no lo procesó.
      \item El SJM muestra la \texttt{GUI\_\allowbreak{}MessageError} indicando que la cuenta pudo haberse vinculado de todos modos y que conviene revisar el correo antes de volver a intentarlo.
      \item Fin del caso de uso.
    \end{cuenum}} \\
   & \cuCelda{\textbf{\mbox{EX-06}: El mensaje recibido está malformado}\par
    \begin{cuenum}[start=1]
      \item El SJM no logra delimitar o deserializar el mensaje conforme al protocolo acordado.
      \item El SJM descarta el mensaje, cierra la conexión y registra en la bitácora local el contenido truncado.
      \item El SJM muestra la \texttt{GUI\_\allowbreak{}MessageError} indicando que la respuesta no pudo interpretarse.
      \item Fin del caso de uso.
    \end{cuenum}} \\
   & \cuCelda{\textbf{\mbox{EX-07}: El servicio de correo no está disponible}\par
    \begin{cuenum}[start=1]
      \item El Servidor de partidas agota los reintentos previstos hacia el Servicio de correo sin obtener confirmación de entrega.
      \item El Servidor de partidas conserva el \textit{Usuario} ya convertido y su \textit{TokenVerificacionCorreo}: la vinculación está hecha y deshacerla devolvería al Invitado a capturarlo todo de nuevo por un fallo que no es suyo (\mbox{RN-07}).
      \item El Servidor de partidas deja \texttt{fecha\_\allowbreak{}ultimo\_\allowbreak{}envio\_\allowbreak{}verificacion} en nulo, de modo que el reenvío del \mbox{FA-09} del \mbox{CU-01} pueda solicitarse de inmediato.
      \item El Servidor de partidas registra el fallo en la bitácora del servidor y responde con \texttt{LINK\_\allowbreak{}ACCOUNT\_\allowbreak{}OK\_\allowbreak{}MAIL\_\allowbreak{}FAILED}.
      \item El SJM despliega la \texttt{GUI\_\allowbreak{}RegistrationSuccess} indicando que la cuenta se vinculó pero el correo de verificación no pudo enviarse, y que puede pedirlo de nuevo al intentar iniciar sesión.
      \item Fin del caso de uso.
    \end{cuenum}} \\ \hline
  \textbf{Postcondiciones} & \cuCelda{\begin{cureglas}
      \item \textbf{\mbox{POST-01}} Si el caso termina por el FN, el \textit{Usuario} conserva el mismo \texttt{id\_\allowbreak{}usuario} y tiene \texttt{tipo\_\allowbreak{}cuenta} en \texttt{REGISTRADA} y \texttt{estado\_\allowbreak{}cuenta} en \texttt{PENDIENTE}.
      \item \textbf{\mbox{POST-02}} Si el caso termina por el FN, todo el progreso anterior sigue intacto: las mismas \textit{EstadisticaModo}, los mismos \textit{UsuarioObjeto}, el mismo \textit{Equipamiento}, las mismas monedas y las mismas \textit{Participacion} en partidas pasadas. No se copió ni se movió ninguna fila.
      \item \textbf{\mbox{POST-03}} Si el caso termina por el FN, existe un registro de \textit{AceptacionTerminos} con la versión aceptada, su idioma, su fecha y la dirección de red de origen.
      \item \textbf{\mbox{POST-04}} Si el caso termina por el FN, existe un \textit{TokenVerificacionCorreo} vigente y se solicitó su envío.
      \item \textbf{\mbox{POST-05}} Si el caso termina por el FN, el \texttt{nickname} definitivo y el \texttt{correo} quedan ocupados y ningún otro \textit{Usuario} puede tomarlos. El \texttt{nickname} provisional de invitado queda libre.
      \item \textbf{\mbox{POST-06}} Si el caso termina por el FN, la \textit{Sesion} de invitado está cerrada y el Invitado no puede seguir jugando hasta verificar su correo con el \mbox{CU-04} e iniciar sesión con el \mbox{CU-01}.
      \item \textbf{\mbox{POST-07}} Si el caso termina por cualquier flujo alterno o por las excepciones \mbox{EX-01} a \mbox{EX-06}, el \textit{Usuario} sigue siendo un invitado con su sesión abierta y su progreso intacto.
    \end{cureglas}} \\
   & \cuCelda{\begin{cureglas}
      \item \textbf{\mbox{POST-08}} Si el caso termina por el \mbox{EX-07}, la cuenta sí quedó vinculada y en \texttt{PENDIENTE}, sin correo enviado y con \texttt{fecha\_\allowbreak{}ultimo\_\allowbreak{}envio\_\allowbreak{}verificacion} en nulo.
    \end{cureglas}} \\ \hline
  \textbf{Reglas de negocio} & \cuCelda{\begin{cureglas}
      \item \textbf{\mbox{RN-01}} La vinculación \textbf{actualiza} el \textit{Usuario} existente. No crea uno nuevo ni copia datos de uno a otro: es lo que hace posible conservar el progreso sin tocar ninguna otra tabla.
      \item \textbf{\mbox{RN-02}} La contraseña no se almacena, transmite ni registra en claro. De ella se guarda únicamente su resumen criptográfico con sal.
      \item \textbf{\mbox{RN-03}} La contraseña tiene al menos ocho caracteres y combina al menos tres de los cuatro grupos de minúsculas, mayúsculas, dígitos y signos. No puede contener el \texttt{nickname} ni el \texttt{correo}.
      \item \textbf{\mbox{RN-04}} La cuenta vinculada nace en \texttt{PENDIENTE}, igual que una creada con el \mbox{CU-02}. Vincular no es verificar: solo el \mbox{CU-04} lleva la cuenta a \texttt{ACTIVA}.
      \item \textbf{\mbox{RN-05}} Un invitado no puede vincularse a un correo ya registrado. Fundir el progreso de dos cuentas obligaría a decidir qué elo, qué monedas y qué historial sobreviven, y no hay respuesta correcta.
      \item \textbf{\mbox{RN-06}} El cambio de tipo de cuenta, la escritura de las credenciales, la \textit{AceptacionTerminos} y el \textit{TokenVerificacionCorreo} ocurren en una sola transacción: o existen todos o no existe ninguno.
      \item \textbf{\mbox{RN-07}} El envío del correo queda fuera de esa transacción. Que el correo falle no puede deshacer una vinculación ya escrita, porque el fallo es del servicio externo y el costo lo pagaría el Invitado.
    \end{cureglas}} \\
   & \cuCelda{\begin{cureglas}
      \item \textbf{\mbox{RN-08}} Vincular cierra la sesión de invitado. La cuenta ya no es de invitado y su forma de entrar pasa a ser el \mbox{CU-01}.
      \item \textbf{\mbox{RN-09}} Todas las validaciones del cliente se repiten en el servidor. Las del cliente existen para dar respuesta inmediata, no para decidir.
      \item \textbf{\mbox{RN-10}} Los datos de una persona que no alcanza la edad mínima no se envían al servidor ni se escriben en ninguna bitácora.
    \end{cureglas}} \\ \hline
  \textbf{Observación} & \cuCelda{\textit{Sobre la cuenta de invitado abandonada.}\par
    Si el Invitado llega hasta el paso 24 y nunca verifica su correo, queda una cuenta en \texttt{PENDIENTE} que ya no puede jugar: perdió la sesión de invitado por el \mbox{RN-08} y no puede iniciar sesión por el \mbox{FA-09} del \mbox{CU-01}. Es el mismo estado en que queda cualquier cuenta creada con el \mbox{CU-02} y sin verificar, así que no introduce un caso nuevo, pero conviene tenerlo presente: vincular es el único momento en que un Invitado puede quedarse sin acceso a un progreso que antes tenía. Por eso el paso 1 advierte de forma visible qué se conserva y qué hace falta para recuperarlo.} \\ \hline
  \textbf{Datos que toca} & \cuCelda{\begin{cudatos}
      \item \textbf{\textit{Sesion}:} lee por token, comprueba que sea de invitado \textit{(paso 8)}
      \item \textbf{\textit{Sesion}:} cierra con \texttt{motivo\_\allowbreak{}cierre} \textit{(paso 23)}
      \item \textbf{\textit{Usuario}:} lee \texttt{nickname} y \texttt{correo} de otros, para comprobar unicidad \textit{(pasos 11, 12)}
      \item \textbf{\textit{Usuario}:} actualiza \texttt{tipo\_\allowbreak{}cuenta}, \texttt{estado\_\allowbreak{}cuenta}, \texttt{nickname}, \texttt{correo}, \texttt{contrasena\_\allowbreak{}hash}, \texttt{contrasena\_\allowbreak{}sal}, \texttt{fecha\_\allowbreak{}nacimiento}, \texttt{fecha\_\allowbreak{}ultimo\_\allowbreak{}envio\_\allowbreak{}verificacion} \textit{(pasos 16, 21)}
      \item \textbf{\textit{PalabraProhibida}:} lee el catálogo con ámbito de nickname \textit{(paso 10)}
      \item \textbf{\textit{AceptacionTerminos}:} crea \textit{(paso 17)}
      \item \textbf{\textit{TokenVerificacionCorreo}:} crea \textit{(paso 18)}
      \item \textbf{\textit{BitacoraAcceso}:} inserta \textit{(paso 22)}
    \end{cudatos}} \\
   & \cuCelda{\begin{cudatos}
      \item \textbf{\textit{EstadisticaModo}, \textit{UsuarioObjeto}, \textit{Equipamiento}, \textit{Participacion}:} \textbf{no se tocan}: conservarlas intactas es el objetivo del caso
    \end{cudatos}} \\ \hline
\end{fichacu}
\clearpage
